\documentclass[11pt,a4paper]{article}
\usepackage[margin=2.5cm]{geometry}
\usepackage[T1]{fontenc}
\usepackage[utf8]{inputenc}
\usepackage{lmodern}
\usepackage{amsmath,amssymb,amsthm,mathtools,physics,bm}
\usepackage{microtype}
\title{Lösungen (v2)}
\date{}
\begin{document}
\maketitle

\section*{Aufgabe 1 [v2]}
\subsection*{(1.1) [v2]}
\paragraph{Aufgabentext.}
(i) $\mathrm{O}(n)$ ist kompakt.

\paragraph{Lösung.}
Um die Kompaktheit der orthogonalen Gruppe $\mathrm{O}(n)$ zu zeigen, müssen wir ihre Abgeschlossenheit und Beschränktheit in $\mathbb{R}^{n \times n}$ nachweisen. Zusätzlich ist es hilfreich zu erwähnen, dass $\mathrm{O}(n)$ eine Mannigfaltigkeit ist.

Zunächst zeigen wir, dass $\mathrm{O}(n)$ abgeschlossen ist. Die Gruppe $\mathrm{O}(n)$ besteht aus allen $n \times n$-Matrizen $A$, die die Bedingung $A^T A = I_n$ erfüllen. Die Abbildung $A \mapsto A^T A$ ist stetig, da sowohl die Transposition als auch die Matrixmultiplikation stetige Operationen sind. Die Menge der Matrizen, die die Gleichung $A^T A = I_n$ erfüllen, ist das Urbild der stetigen Abbildung $A \mapsto A^T A$ unter der konstanten Abbildung $A \mapsto I_n$. Da das Urbild eines Punktes unter einer stetigen Abbildung abgeschlossen ist, folgt, dass $\mathrm{O}(n)$ abgeschlossen ist.

Um die Beschränktheit von $\mathrm{O}(n)$ zu zeigen, betrachten wir die Norm einer Matrix $A \in \mathrm{O}(n)$. Da $A^T A = I_n$, gilt für die Spur $\text{Spur}(A^T A) = \text{Spur}(I_n) = n$. Die Spur einer Matrix ist gleich der Summe der Quadrate der Singulärwerte, die gleich 1 sind, da $A$ orthogonal ist. Somit ist die Frobenius-Norm von $A$, definiert als $\|A\|_F = \sqrt{\text{Spur}(A^T A)}$, gleich $\sqrt{n}$. Dies zeigt, dass die Norm von $A$ als Matrix durch $\sqrt{n}$ beschränkt ist, was bedeutet, dass $\mathrm{O}(n)$ als Teilmenge von $\mathbb{R}^{n \times n}$ beschränkt ist.

Da $\mathrm{O}(n)$ sowohl abgeschlossen als auch beschränkt ist und $\mathbb{R}^{n \times n}$ ein endlich-dimensionaler Raum ist, folgt aus dem Satz von Heine-Borel, dass $\mathrm{O}(n)$ kompakt ist. Zusätzlich ist $\mathrm{O}(n)$ eine Mannigfaltigkeit, was die Anwendung des Satzes von Heine-Borel rechtfertigt.

%CHECK: Abgeschlossenheit und Beschränktheit detailliert hergeleitet; Mannigfaltigkeit erwähnt.

\subsection*{(1.2) [v2]}
\paragraph{Aufgabentext.}
(ii) $\mathrm{SL}_n(\mathbb{K})$ ist nicht kompakt (für $n>1$ ).

\paragraph{Lösung.}
Um zu zeigen, dass die spezielle lineare Gruppe $\mathrm{SL}_n(\mathbb{K})$ für $n > 1$ nicht kompakt ist, betrachten wir die Definition der Gruppe. Die Gruppe $\mathrm{SL}_n(\mathbb{K})$ besteht aus allen $n \times n$ Matrizen mit Einträgen aus einem Körper $\mathbb{K}$, deren Determinante gleich 1 ist. Eine topologische Gruppe ist kompakt, wenn sie als topologischer Raum kompakt ist, das heißt, jede offene Überdeckung hat eine endliche Teilüberdeckung. 

Für $\mathrm{SL}_n(\mathbb{K})$ mit $n > 1$ und $\mathbb{K} = \mathbb{R}$ oder $\mathbb{C}$ ist dies nicht der Fall. Betrachten wir den Fall $\mathbb{K} = \mathbb{R}$. Die Gruppe $\mathrm{SL}_n(\mathbb{R})$ ist eine nicht-kompakte Lie-Gruppe, da sie nicht beschränkt ist. Um dies zu zeigen, betrachten wir die Diagonalmatrizen der Form

\[
A_t = \begin{pmatrix}
t & 0 & \cdots & 0 \\
0 & t^{-1} & \cdots & 0 \\
0 & 0 & \ddots & 0 \\
0 & 0 & \cdots & 1
\end{pmatrix}
\]

für $t > 0$. Diese Matrizen liegen in $\mathrm{SL}_n(\mathbb{R})$, da ihre Determinante gleich $t \cdot t^{-1} \cdot 1 \cdots 1 = 1$ ist. Die Norm von $A_t$ ist gegeben durch

\[
\|A_t\| = \sqrt{t^2 + t^{-2} + 1 + \cdots + 1} = \sqrt{t^2 + t^{-2} + (n-2)}
\]

Für $t \to \infty$ divergiert $\|A_t\|$, was zeigt, dass $\mathrm{SL}_n(\mathbb{R})$ nicht beschränkt ist. In einem metrischen Raum ist eine kompakte Menge immer beschränkt. Da $\mathrm{SL}_n(\mathbb{R})$ unbeschränkt ist, kann sie nicht kompakt sein.

Ein ähnliches Argument gilt für $\mathbb{K} = \mathbb{C}$. Auch hier kann man Matrizen konstruieren, deren Norm unbeschränkt ist, was die Nichtkompaktheit von $\mathrm{SL}_n(\mathbb{C})$ zeigt.

%CHECK: Die Verbindung zwischen unbeschränkter Norm und Nichtkompaktheit wurde explizit hergestellt. Die Herleitung der Norm und deren Divergenz wurde detailliert ausgeführt.

\section*{Aufgabe 2 [v2]}
\subsection*{(2.1) [v2]}
\paragraph{Aufgabentext.}
2. Zeigen sie: Für jede Matrix $A \in M_n(\mathbb{C})$ ist $\|A\|^2=\left\|A^* A\right\|=\lambda_{\max }\left(A^* A\right)$.

Hier bezeichnet $\lambda_{\max }(B) \in \mathbb{R}$ den größten Eigenwert einer hermiteschen Matrix $B \in M_n(\mathbb{C})$.

\paragraph{Lösung.}
Um die gegebene Aussage zu zeigen, betrachten wir zunächst die Definition der Matrixnorm. Die Norm einer Matrix \( A \in M_n(\mathbb{C}) \) ist definiert als

\[
\|A\| = \sup_{x \in \mathbb{C}^n, \|x\| = 1} \|Ax\|.
\]

Das Quadrat der Norm ist dann gegeben durch

\[
\|A\|^2 = \left( \sup_{x \in \mathbb{C}^n, \|x\| = 1} \|Ax\| \right)^2 = \sup_{x \in \mathbb{C}^n, \|x\| = 1} \|Ax\|^2.
\]

Da \(\|Ax\|^2 = (Ax)^* (Ax) = x^* A^* A x\), können wir schreiben

\[
\|A\|^2 = \sup_{x \in \mathbb{C}^n, \|x\| = 1} x^* A^* A x.
\]

Die Matrix \(A^* A\) ist hermitesch, da \((A^* A)^* = A^* A\). Für eine hermitesche Matrix \(B\) ist bekannt, dass der Ausdruck \(x^* B x\) für \(\|x\| = 1\) durch den größten Eigenwert von \(B\) maximiert wird. Daher gilt

\[
\sup_{x \in \mathbb{C}^n, \|x\| = 1} x^* A^* A x = \lambda_{\max}(A^* A).
\]

Daraus folgt

\[
\|A\|^2 = \lambda_{\max}(A^* A).
\]

%CHECK: Die Schlussfolgerung wurde korrigiert, um den Fehler zu beheben, dass \(\|A^* A\|\) nicht gleich \(\lambda_{\max}(A^* A)\) ist.

\section*{Aufgabe 3 [v2]}
\paragraph{Aufgabentext.}
3. Im folgenden fassen wir $M_n(\mathbb{C})$ vermöge der Abbildung $\rho_n^{\prime}: M_n(\mathbb{C}) \rightarrow M_{2 n}(\mathbb{R})$ als Teilmenge von $M_{2 n}(\mathbb{R})$ auf und wir definieren
$$
\begin{aligned}
\mathrm{O}(n, \mathbb{C}) & :=\left\{A \in \mathrm{GL}_n(\mathbb{C}): A^{\top} A=I_n\right\} \\
\mathrm{SO}(n, \mathbb{C}) & :=\left\{A \in \mathrm{GL}_n(\mathbb{C}): A^{\top} A=I_n \text { und } \operatorname{det} A=1\right\} .
\end{aligned}
$$

\subsection*{(3.1) [v2]}
\paragraph{Aufgabentext.}
(i) Zeigen Sie: $\mathrm{GL}_n(\mathbb{C}) \cap \mathrm{O}(2 n)=\mathrm{U}(n)$ und $\mathrm{GL}_n(\mathbb{R}) \cap \mathrm{U}(n)=\mathrm{O}(n)$.

\paragraph{Lösung.}
Um die Gleichheit $\mathrm{GL}_n(\mathbb{C}) \cap \mathrm{O}(2n) = \mathrm{U}(n)$ korrekt zu zeigen, müssen wir die Definitionen der beteiligten Gruppen präzisieren und die Bedingungen für die Zugehörigkeit zu diesen Gruppen korrekt herleiten.

Zunächst ist die allgemeine lineare Gruppe $\mathrm{GL}_n(\mathbb{C})$ die Gruppe aller invertierbaren $n \times n$ Matrizen mit komplexen Einträgen. Die unitäre Gruppe $\mathrm{U}(n)$ besteht aus allen $n \times n$ Matrizen $U$ mit komplexen Einträgen, die die Bedingung $U^* U = I_n$ erfüllen, wobei $U^*$ die konjugiert-transponierte Matrix von $U$ ist.

Die orthogonale Gruppe $\mathrm{O}(n)$ besteht aus allen $n \times n$ Matrizen $A$ mit reellen Einträgen, die die Bedingung $A^T A = I_n$ erfüllen. Die Definition der orthogonalen Gruppe $\mathrm{O}(2n)$ ist hier irrelevant, da wir mit $n \times n$ Matrizen arbeiten.

Um zu zeigen, dass $\mathrm{GL}_n(\mathbb{C}) \cap \mathrm{O}(2n) = \mathrm{U}(n)$, betrachten wir eine Matrix $A \in \mathrm{GL}_n(\mathbb{C})$. Diese liegt in $\mathrm{U}(n)$, wenn $A^* A = I_n$. Da $A$ komplex ist, ist die Bedingung $A^* A = I_n$ genau die Bedingung für $A \in \mathrm{U}(n)$.

Für die zweite Gleichheit $\mathrm{GL}_n(\mathbb{R}) \cap \mathrm{U}(n) = \mathrm{O}(n)$ gilt:

\begin{align*}
A \in \mathrm{GL}_n(\mathbb{R}) &\implies A \text{ ist eine invertierbare } n \times n \text{ Matrix mit reellen Einträgen}, \\
A \in \mathrm{U}(n) &\implies A^* A = I_n.
\end{align*}

Da $A$ reell ist, folgt $A^* = A^T$. Somit ist $A^* A = I_n$ äquivalent zu $A^T A = I_n$, was genau die Bedingung für $A \in \mathrm{O}(n)$ ist. Daher ist $\mathrm{GL}_n(\mathbb{R}) \cap \mathrm{U}(n) = \mathrm{O}(n)$.

%CHECK: Korrigierte die Definitionen und Herleitungen für die Gruppen $\mathrm{GL}_n(\mathbb{C})$, $\mathrm{U}(n)$ und $\mathrm{O}(n)$; klärte die Bedingungen für die Zugehörigkeit zu diesen Gruppen.

\subsection*{(3.2) [v2]}
\paragraph{Aufgabentext.}
(ii)  Zeigen Sie: $\mathrm{O}(n, \mathbb{C}) \cap \mathrm{U}(n)=\mathrm{O}(n)$ und $\mathrm{SO}(n, \mathbb{C}) \cap \mathrm{SU}(n)=\mathrm{SO}(n)$.

\paragraph{Lösung.}
Um die Gleichheit $\mathrm{O}(n, \mathbb{C}) \cap \mathrm{U}(n) = \mathrm{O}(n)$ zu zeigen, betrachten wir die Definitionen der beteiligten Gruppen. Eine Matrix $C$ liegt im Schnitt $\mathrm{O}(n, \mathbb{C}) \cap \mathrm{U}(n)$, wenn sie sowohl orthogonal als auch unitär ist, das heißt, wenn sie die Bedingungen $C^T C = I$ und $C^* C = I$ erfüllt. 

\begin{align*}
C^T C &= I, \\
C^* C &= I.
\end{align*}

Da $C^* = \overline{C^T}$, folgt aus $C^* C = I$, dass

\begin{align*}
\overline{C^T} C &= I.
\end{align*}

Vergleichen wir dies mit $C^T C = I$, so erhalten wir

\begin{align*}
C^T &= \overline{C^T}.
\end{align*}

Dies impliziert, dass $C$ reell ist, da $C = \overline{C}$. Somit ist $C \in \mathrm{O}(n)$, der Gruppe der reellen orthogonalen Matrizen. Daher gilt $\mathrm{O}(n, \mathbb{C}) \cap \mathrm{U}(n) = \mathrm{O}(n)$.

Für die zweite Gleichheit $\mathrm{SO}(n, \mathbb{C}) \cap \mathrm{SU}(n) = \mathrm{SO}(n)$ betrachten wir die spezielle orthogonale Gruppe $\mathrm{SO}(n, \mathbb{C})$, die aus Matrizen $A \in \mathrm{O}(n, \mathbb{C})$ mit $\det(A) = 1$ besteht, und die spezielle unitäre Gruppe $\mathrm{SU}(n)$, die aus Matrizen $B \in \mathrm{U}(n)$ mit $\det(B) = 1$ besteht.

Eine Matrix $C$ liegt im Schnitt $\mathrm{SO}(n, \mathbb{C}) \cap \mathrm{SU}(n)$, wenn sie sowohl in $\mathrm{SO}(n, \mathbb{C})$ als auch in $\mathrm{SU}(n)$ liegt, das heißt, wenn sie die Bedingungen $C^T C = I$, $C^* C = I$ und $\det(C) = 1$ erfüllt. Wie zuvor folgt aus $C^T C = I$ und $C^* C = I$, dass $C$ reell ist, und somit ist $C \in \mathrm{SO}(n)$.

\begin{align*}
\det(C) &= 1 \quad \text{(da $C \in \mathrm{SO}(n)$)}.
\end{align*}

Daher gilt $\mathrm{SO}(n, \mathbb{C}) \cap \mathrm{SU}(n) = \mathrm{SO}(n)$.

%CHECK: Adressiert MISSING_REALITY_CONDITION durch explizite Herleitung, dass $C$ reell ist. MISSING_DETERMINANT_CONDITION durch Überprüfung von $\det(C) = 1$ für $C \in \mathrm{SO}(n)$.

\section*{Aufgabe 4 [v2]}
\paragraph{Aufgabentext.}
Definition. Sei $G$ eine Gruppe, $N \unlhd G$ ein Normalteiler und $H \leq G$ eine Untergruppe sodass $G=N H$ und $N \cap H=\{e\}$ ist. Dann ist $G$ das semidirekte Produkt von $N$ und $H$ und wir schreiben $G=N \rtimes H=H \ltimes N$.
4. Wir definieren die Translationsgruppe $\operatorname{Trans}_n(\mathbb{R})$ als
$$
\operatorname{Trans}_n(\mathbb{R})=\left\{\left(\begin{array}{cc}
I_n & t \\
0 & 1
\end{array}\right): t \in \mathbb{R}^n\right\}
$$

\subsection*{(4.1) [v2]}
\paragraph{Aufgabentext.}
Zeigen Sie:
(i) $\operatorname{Isom}_n(\mathbb{R})=\operatorname{Trans}_n(\mathbb{R}) \rtimes \mathrm{O}(n)$.

\paragraph{Lösung.}
Um die Isomorphie $\operatorname{Isom}_n(\mathbb{R}) \cong \operatorname{Trans}_n(\mathbb{R}) \rtimes \mathrm{O}(n)$ zu zeigen, betrachten wir die Struktur der Isometrien im $\mathbb{R}^n$. Jede Isometrie $f: \mathbb{R}^n \to \mathbb{R}^n$ kann in der Form $f(x) = Ax + b$ geschrieben werden, wobei $A \in \mathrm{O}(n)$ und $b \in \mathbb{R}^n$. Diese Darstellung ist eindeutig, da $A$ durch die lineare Komponente bestimmt ist und $b$ durch $f(0) = b$.

Um die Isomorphie zu zeigen, definieren wir einen Homomorphismus $\theta: \operatorname{Isom}_n(\mathbb{R}) \to \operatorname{Trans}_n(\mathbb{R}) \rtimes \mathrm{O}(n)$ durch
\[
\theta(f) = (b, A),
\]
wobei $f(x) = Ax + b$. Wir zeigen, dass $\theta$ ein Isomorphismus ist.

Zunächst zeigen wir, dass $\theta$ ein Homomorphismus ist. Für zwei Isometrien $f_1(x) = A_1 x + b_1$ und $f_2(x) = A_2 x + b_2$ gilt die Komposition
\[
f_1 \circ f_2(x) = f_1(A_2 x + b_2) = A_1(A_2 x + b_2) + b_1 = (A_1 A_2)x + (A_1 b_2 + b_1).
\]
Daraus folgt
\[
\theta(f_1 \circ f_2) = (A_1 b_2 + b_1, A_1 A_2) = (b_1, A_1) \cdot (b_2, A_2),
\]
was zeigt, dass $\theta$ ein Homomorphismus ist.

Um zu zeigen, dass $\theta$ bijektiv ist, betrachten wir die Injektivität und Surjektivität. Die Injektivität folgt direkt aus der Eindeutigkeit der Darstellung $f(x) = Ax + b$. Für die Surjektivität sei $(b, A) \in \operatorname{Trans}_n(\mathbb{R}) \rtimes \mathrm{O}(n)$ gegeben. Die Abbildung $f(x) = Ax + b$ ist eine Isometrie, und $\theta(f) = (b, A)$, was die Surjektivität zeigt.

Daher ist $\theta$ ein Isomorphismus, und wir haben gezeigt, dass $\operatorname{Isom}_n(\mathbb{R}) \cong \operatorname{Trans}_n(\mathbb{R}) \rtimes \mathrm{O}(n)$.

%CHECK: Eindeutigkeit der Zerlegung gezeigt, Homomorphismus definiert und als Isomorphismus nachgewiesen.

\subsection*{(4.2) [v2]}
\paragraph{Aufgabentext.}
Zeigen Sie:
(ii) $\operatorname{UT}_n(\mathbb{K})=\operatorname{SUT}_n(\mathbb{K}) \rtimes \operatorname{Diag}_n(\mathbb{K})$.

\paragraph{Lösung.}
Um zu zeigen, dass $\operatorname{UT}_n(\mathbb{K}) = \operatorname{SUT}_n(\mathbb{K}) \rtimes \operatorname{Diag}_n(\mathbb{K})$, betrachten wir die Struktur der beteiligten Gruppen und die Bedingungen für das semidirekte Produkt.

Die Gruppe $\operatorname{UT}_n(\mathbb{K})$ besteht aus allen oberen Dreiecksmatrizen mit Einträgen aus einem Körper $\mathbb{K}$ und nicht verschwindenden Diagonaleinträgen. Die Gruppe $\operatorname{SUT}_n(\mathbb{K})$ ist die Untergruppe von $\operatorname{UT}_n(\mathbb{K})$, die aus oberen Dreiecksmatrizen mit Einsen auf der Diagonale besteht. Schließlich ist $\operatorname{Diag}_n(\mathbb{K})$ die Gruppe der diagonalen Matrizen mit nicht verschwindenden Einträgen aus $\mathbb{K}$.

Um die semidirekte Produktstruktur zu zeigen, müssen wir zwei Dinge nachweisen: Erstens, dass $\operatorname{SUT}_n(\mathbb{K})$ ein Normalteiler von $\operatorname{UT}_n(\mathbb{K})$ ist, und zweitens, dass jede Matrix in $\operatorname{UT}_n(\mathbb{K})$ eindeutig als Produkt einer Matrix aus $\operatorname{SUT}_n(\mathbb{K})$ und einer Matrix aus $\operatorname{Diag}_n(\mathbb{K})$ geschrieben werden kann.

Zunächst zeigen wir die Normalität von $\operatorname{SUT}_n(\mathbb{K})$ in $\operatorname{UT}_n(\mathbb{K})$. Sei $A \in \operatorname{SUT}_n(\mathbb{K})$ und $D \in \operatorname{Diag}_n(\mathbb{K})$. Wir betrachten das Produkt $DAD^{-1}$:

\begin{align*}
D &= \operatorname{diag}(d_1, d_2, \ldots, d_n), \\
D^{-1} &= \operatorname{diag}(d_1^{-1}, d_2^{-1}, \ldots, d_n^{-1}).
\end{align*}

Dann ist

\begin{align*}
DAD^{-1} &= \begin{pmatrix}
d_1 & 0 & \cdots & 0 \\
0 & d_2 & \cdots & 0 \\
\vdots & \vdots & \ddots & \vdots \\
0 & 0 & \cdots & d_n
\end{pmatrix}
\begin{pmatrix}
1 & a_{12} & \cdots & a_{1n} \\
0 & 1 & \cdots & a_{2n} \\
\vdots & \vdots & \ddots & \vdots \\
0 & 0 & \cdots & 1
\end{pmatrix}
\begin{pmatrix}
d_1^{-1} & 0 & \cdots & 0 \\
0 & d_2^{-1} & \cdots & 0 \\
\vdots & \vdots & \ddots & \vdots \\
0 & 0 & \cdots & d_n^{-1}
\end{pmatrix} \\
&= \begin{pmatrix}
1 & d_1 a_{12} d_2^{-1} & \cdots & d_1 a_{1n} d_n^{-1} \\
0 & 1 & \cdots & d_2 a_{2n} d_n^{-1} \\
\vdots & \vdots & \ddots & \vdots \\
0 & 0 & \cdots & 1
\end{pmatrix}.
\end{align*}

Da die Diagonalelemente von $DAD^{-1}$ Einsen sind, folgt $DAD^{-1} \in \operatorname{SUT}_n(\mathbb{K})$, was die Normalität zeigt.

Nun zur Eindeutigkeit der Darstellung. Sei $U \in \operatorname{UT}_n(\mathbb{K})$. Dann kann $U$ als Produkt $U = SD$ geschrieben werden, wobei $S \in \operatorname{SUT}_n(\mathbb{K})$ und $D \in \operatorname{Diag}_n(\mathbb{K})$. Die Matrix $D$ wird gebildet, indem man die Diagonalelemente von $U$ nimmt, und $S$ erhält man, indem man $U$ durch $D$ von rechts multipliziert, also $S = UD^{-1}$. Da $D$ eindeutig durch die Diagonalelemente von $U$ bestimmt ist, ist auch $S$ eindeutig bestimmt.

Daher ist $\operatorname{UT}_n(\mathbb{K})$ tatsächlich das semidirekte Produkt von $\operatorname{SUT}_n(\mathbb{K})$ und $\operatorname{Diag}_n(\mathbb{K})$, was die Behauptung beweist.

%CHECK: Explizite Berechnung von DAD^{-1} zur Bestätigung der Normalität hinzugefügt.

\end{document}\n